\section{The ratchet, III: set-up and linear algebra}\label{sec:setup}

\subsection{Set-up}

Set-up is the largest term on the Koblitz pipelines (\S\ref{sec:model}).
\begin{itemize}
\item \emph{Stacked $n=83$ producer} (\texttt{cryptanalysis} \#938). This
  stacks the cryptanalysis repository's set-up optimisations into one
  producer. It made end-to-end cold $3.2391\times$ $[2.9809,3.5240]$
  faster than its own baseline. Paired within that run, the cold gap to
  rho fell from $305.8$ to $94.4$. The online ratio moved from $0.133$ to
  $0.139$, a small online regression. An integrity audit upheld the claim
  within a narrower scope (\art{crypto-autoresearcher:research/sheriff/SHF-20261016-3caa6f/findings.yaml},
  C17--C19):
  \begin{itemize}
  \item it is one run on one runner class, with no replay on another host;
  \item about $1.48\times$ of the gain is host-independent operation
    counts, and the rest is wall time;
  \item factor-base construction is charged to neither arm;
  \item the repository's ``$360\to94.4$'' wording and a set-up field
    compared across rows with different meanings were judged
    overclaimed.
  \end{itemize}
  We report the scoped figure.
\item \emph{Cache-blocked pair table.} Building the prime-field pair table
  in cache-sized blocks cuts set-up to $0.24$ $[0.22,0.26]$ of the batched
  build and $0.13$ $[0.12,0.14]$ of the scalar build at $N=16{,}384$. That
  is $28.2$\,s instead of $117$\,s and $224$\,s, and peak memory falls from
  $9.15$\,GiB to $4.65$\,GiB. All 456 of 456 executions are identical, and
  the A/A ratio is $1.02$ $[0.70,1.46]$
  (\art{research/ic_oracle_setup_blocked_20261020/}). Two caveats apply.
  The host earned L1, so the wall ratios are descriptive. And counted $S$
  is unchanged, so this is a wall-time gain that leaves the gap in the
  counted unit as it was (rho is $32.96$, $74.31$ and $269.03\times$
  cheaper at $N=1{,}024$ to $16{,}384$, $w=0$).
\item \emph{u64 root batching.} Batching root extraction in 64-bit lanes
  gives $3.26\times$ whole-process. The rank-stage and per-target cells were
  corrected on 20 October from $2.95\times$ and $6.09\times$ to $2.47\times$
  and $5.00\times$; the originals had compared the candidate's fastest run
  with the baseline median (benchmarks row \texttt{20261015-u64-rootbatch}).
\end{itemize}

\subsection{Linear algebra}

\begin{itemize}
\item \emph{CADO-style filtering and Wiedemann.} On synthetic $m=3$
  matrices over a 62-bit prime, a filter followed by
  Wiedemann~\cite{wiedemann1986} beats the legacy solver by $2.48\times$
  at $2^{14}$ columns and $10.40\times$ $[9.99,11.06]$ at $2^{18}$. At
  $2^{18}$ it is $1.93\times$ $[1.74,1.96]$ faster than CADO-NFS itself
  (\art{research/la_engines_20261014/}).
\item \emph{Inside ecbench.} Wired into the end-to-end pipeline, the same
  solver leaves counted $S$ at $1.000$ and cuts linear-algebra wall time
  to $0.25$--$0.72$. Linear algebra is only 0.1--1.9\% of these runs, and
  the best cell is still $149\times$ rho
  (\art{research/ic_linalg_pipeline_20261019/}).
\item \emph{Deferred sparse solvers.} A negative: no deferred sparse
  solver lowers end-to-end $S$ ($1.206\times$ on the $k=4$ ladder;
  $5.72\to1.030\times$ on the scaling curve). These solvers answer only at
  full rank, which is a completion-rule effect, not a linear-algebra one
  (\art{research/sparse_la_timing_20261018/README.md}).
\item \emph{Sparse FGLM in the Jacobian-cover pipeline} cuts $S$ by
  $1.28$--$1.29\times$. For example, IC/rho falls from $851$ to $664$ at
  $p=53$ (\art{research/jv_cover_sparse_fglm_20261018/}). Both are lower
  bounds.
\end{itemize}

\paragraph{Reading.} The set-up results are the largest end-to-end
constant-factor gains of the period: $3.24\times$ at $n=83$, and $4$--$8\times$
in prime set-up wall. Linear algebra is a large stage gain on a term that
is small on the curves measured. It becomes the term that matters on the
genus-three family (\S\ref{sec:families}), where large primes and
structured Gaussian elimination moved the exponent.
